
\section{Sīl'uddesa-pāṭha}
\label{siluddesa-patha}

\begin{intro}
	\anglebracketleft\ \hspace{-0.5mm}Handa mayaṁ sīl'uddesa-pāṭhaṁ bhaṇāmase \hspace{-0.5mm}\anglebracketright\
\end{intro}

%% \medskip

\begin{hanging-indent}
Bhāsitam'idaṁ tena bhagavatā jānatā passatā arahatā sammāsambuddhena:
\end{hanging-indent}
\begin{hanging-indent-together}
``Sampanna-sīlā bhikkhave viharatha sampanna-pāṭimokkhā.
\end{hanging-indent-together}
\begin{hanging-indent-together}
Pāṭimokkha-saṁvara-saṁvutā viharatha ācāra-gocara-sampannā.
\end{hanging-indent-together}
\begin{hanging-indent-together}
Aṇumattesu vajjesu bhaya-dassāvī samādāya sikkhatha sikkhāpadesū'ti.
\end{hanging-indent-together}
\begin{hanging-indent-together}
Tasmā't'ih'amhehi sikkhitabbaṁ.
\end{hanging-indent-together}

\begin{hanging-indent}
``Sampanna-sīlā viharissāma sampanna-pāṭimokkhā.
\end{hanging-indent}
\begin{hanging-indent-together}
Pāṭimokkha-saṁvara-saṁvutā viharissāma ācāra-gocara-sampannā.
\end{hanging-indent-together}
\begin{hanging-indent-together}
Aṇumattesu vajjesu bhaya-dassāvī samādāya sikkhissāma sikkhāpadesū'ti.
\end{hanging-indent-together}
\begin{hanging-indent-together}
Evañ'hi no sikkhitabbaṁ.
\end{hanging-indent-together}

\suttaRef{[MN 6]}

\clearpage



\section{The Summary on Virtue}
\label{siluddesa-patha-english}

\begin{intro}
	\anglebracketleft\ \hspace{-0.5mm}Let us recite the summary on virtue: \hspace{-0.5mm}\anglebracketright\
\end{intro}

\begin{hanging-indent}
This was said by the Blessed One, the One who Knows, the One who Sees, the Worthy One, the Perfectly Enlightened One:
\end{hanging-indent}
\begin{hanging-indent-together}
``Live consummate in virtue, bhikkhus, and consummate in the Pāṭimokkhā.
\end{hanging-indent-together}
\begin{hanging-indent-together}
Live restrained with the restrain of the Pāṭimokkhā, consummate in your behaviour and sphere of activity.
\end{hanging-indent-together}
\begin{hanging-indent-together}
Train yourselves, having undertaken the training rules, seeing danger in the slightest faults.''
\end{hanging-indent-together}

\begin{hanging-indent}
Therefore we should train ourselves:
\end{hanging-indent}
\begin{hanging-indent-together}
``We will live consummate in virtue, consummate in the Pāṭimokkhā.
\end{hanging-indent-together}
\begin{hanging-indent-together}
We will live restrained with the restrain of the Pāṭimokkhā, consummate in our behaviour and sphere of activity.
\end{hanging-indent-together}
\begin{hanging-indent-together}
We will train ourselves, having undertaken the training rules, seeing danger in the slightest faults.''
\end{hanging-indent-together}
\begin{hanging-indent-together}
That's how we should train ourselves.
\end{hanging-indent-together}

\suttaRef{[MN 6]}

\clearpage

